% First part belongs immediately after the three-state inverse theorem.
\subsection{Conditioning of the three-state inverse}\label{subsec:conditioning}

Local Lipschitz inversion is a pointwise statement. To quantify its constants,
write $\theta=(r,s,a,c,b,d)$ and let $\phi(\theta)\in\mathbb R^{12}$ collect
the joint transform entries at three distinct positive arguments. With
$K_\lambda=(\lambda I-T)^{-1}$, its full six-rate derivative is
\begin{equation}
 \partial_j\widehat\Psi(\lambda)
 =RK_\lambda(\partial_jT)K_\lambda B
   +RK_\lambda(\partial_jB).\label{eq:transform-jacobian}
\end{equation}
The second term is essential when the observed rates $r,s$ are unknown.
Absolute-rate errors are governed by $J_\theta=D\phi$, whereas relative-rate
errors are governed by $J_{\log}=J_\theta\operatorname{diag}(\theta)$.
These describe different inverse questions. Raw singular values carry time
units; logarithmic coordinates are dimensionless in the fixed time convention.

The rational inverse in Theorem~\ref{thm:inverse} has a derivative that is a
left inverse of $J_\theta$. This remains true at an irreducible tree by
continuity from the positive interior. Consequently the raw derivative has
rank six there. Sparse limits need not make absolute-rate recovery ill
conditioned. For example, at the unit tree and $\lambda=(0.2,1,5)$, its smallest
singular value is approximately $0.00744$.

For the family~\eqref{eq:treefamily}, fix $0<e<1/2$ and let $k\downarrow0$.
The rational formulas also remain nonsingular at this irreducible one-way
limit, although the limiting model itself is outside our reciprocal class.
Thus $J_\theta$ converges to a full-rank matrix while
\begin{equation}
 s_{\min}(J_{\log})=\Theta(k),\qquad
 \operatorname{cond}(J_{\log})=\Theta(k^{-1}).\label{eq:log-conditioning}
\end{equation}
Indeed $s_{\min}(J_\theta D)\ge s_{\min}(J_\theta)\min_jD_{jj}$ for
$D=\operatorname{diag}(\theta)$, and testing the unit $b$ coordinate gives
the matching upper bound. The other five rate coordinates stay positive,
so the limiting nullspace of $J_{\log}$ is exactly the $b=k$ direction.
Its weakest right singular vector converges to that direction, up to sign.
Meanwhile $\partial_b\sig=-J_0(e)/b+O(\log(1/b))$, with
$J_0(e)=e(1-e)/(3+e-e^2)>0$. Relative-rate and entropy sensitivity therefore
diverge even where absolute-rate inversion remains regular.

The choice of Laplace arguments also matters. As three distinct arguments
coalesce, the limiting repeated transform matrix has at most four independent
entries for six rates, so a singular value tends to zero. Arguments all tending
to zero or infinity similarly lose informative variation. Hence ``any three
distinct positive points'' guarantees exact recoverability, with no uniform
conditioning guarantee over designs or sources.

\begin{figure}[t]
 \centering
 \includegraphics[width=\textwidth]{figures/conditioning.pdf}
 \caption{Conditioning of all six unknown three-state rates. (a) Log-rate
 condition number for $(\lambda_1,\lambda_2,\lambda_3)=(c/v,c,cv)$ at
 $(r,s,a,c,b,d)=(1.2,0.8,0.9,0.4,0.3,1.1)$; the horizontal $c$ denotes the
 Laplace center, not the rate coordinate. Nearly coincident points and
 inappropriate scales degrade inversion. (b) At fixed $e=0.2$ in
 Eq.~\eqref{eq:treefamily}, the raw minimum singular value stays positive as
 $k$ decreases, whereas its log-rate counterpart scales as $k$.
 (c) The weakest log-rate direction tends to the reverse-rate coordinate
 $b=k$. Panels (b,c) use $\lambda=(0.2,1,5)$. These unweighted Euclidean
 diagnostics are not Fisher information; statistical noise must also be
 considered.}\label{fig:conditioning}
\end{figure}

We compared all 165 triples from
$\{0.025,0.05,0.1,0.2,0.5,1,2,5,10,20,40\}$ at the reference in
Fig.~\ref{fig:conditioning}. The largest minimum log singular value occurs
at $(0.2,0.5,5)$. The smallest oracle asymptotic entropy standard error for
the moment estimator below instead occurs at $(0.025,10,20)$: its
$\mathrm{SE}\sqrt n$ is $1.176$, compared with $1.238$ for $(0.2,1,5)$.
The better entropy design has a worse Euclidean condition number.
These are finite-grid comparisons at one reference, not universal optima.
A practical starting point is a separated logarithmic triple spanning a
pilot waiting-time scale, followed by a covariance-aware assessment of the
target uncertainty. Design selection from the same data needs its own
selection correction or an independent pilot.

